%% 自動生成（tools/mknext.py）。各章の道しるべ「この章の先に」から、後ろの巻の章を集めたもの
\begin{longtblr}{colspec={Q[l,m,wd=58mm] X[l,m]}, width=\linewidth, rowhead=1, row{1}={bg=bkVol!14, font=\sffamily\bfseries\footnotesize},
  hlines={0.4pt,bkLine}, rowsep=1.5pt, cells={font=\footnotesize}}
この\R{巻}{かん}の\R{章}{しょう} & この\R{道具}{どう|ぐ}を\R{使}{つか}う\R{章}{しょう} \\
{\bknextpill{3-1}\ \R{変}{か}わる\R{量}{りょう}と\R{変}{か}わらない\R{量}{りょう}} & {\R{第}{だい}4\R{巻}{かん}から\R{先}{さき}：\bknextpill{5-4}\,\bknextpill{5-5}} \\
{\bknextpill{3-2}\ \R{場所}{ば|しょ}を\R{数}{かず}で\R{表}{あらわ}す：\R{座標}{ざ|ひょう}} & {\R{第}{だい}4\R{巻}{かん}から\R{先}{さき}：\bknextpill{4-4}\,\bknextpill{4-10}\,\bknextpill{4-13}\,\bknextpill{4-15}} \\
{\bknextpill{3-3}\ \R{関数}{かん|すう}は「\R{入}{い}れると\R{出}{で}てくる」ルール} & {\R{第}{だい}4\R{巻}{かん}から\R{先}{さき}：\bknextpill{5-1}\,\bknextpill{5-3}\,\bknextpill{5-6}\,\bknextpill{9-8}} \\
{\bknextpill{3-4}\ グラフをかこう・\R{読}{よ}もう} & {\R{第}{だい}4\R{巻}{かん}から\R{先}{さき}：\bknextpill{5-1}\,\bknextpill{5-9}} \\
{\bknextpill{3-5}\ グラフの\R{傾}{かたむ}きは「\R{変化}{へん|か}の\R{速}{はや}さ」} & {\R{第}{だい}4\R{巻}{かん}から\R{先}{さき}：\bknextpill{5-1}\,\bknextpill{5-3}\,\bknextpill{8-2}} \\
{\bknextpill{3-6}\ 2\R{本}{ほん}の\R{直線}{ちょく|せん}が\R{交}{まじ}わるところ：\R{連立方程式}{れん|りつ|ほう|てい|しき}} & {\R{第}{だい}4\R{巻}{かん}から\R{先}{さき}：\bknextpill{7-1}\,\bknextpill{9-10}} \\
{\bknextpill{3-7}\ \R{反比例}{はん|ぴ|れい}のグラフ} & {\R{第}{だい}4\R{巻}{かん}から\R{先}{さき}：\bknextpill{5-4}\,\bknextpill{5-14}} \\
{\bknextpill{3-8}\ かっこをはずす：\R{展開}{てん|かい}と\R{平方完成}{へい|ほう|かん|せい}} & {\R{第}{だい}4\R{巻}{かん}から\R{先}{さき}：\bknextpill{5-4}\,\bknextpill{9-7}} \\
{\bknextpill{3-9}\ \R{放物線}{ほう|ぶつ|せん}と $y=ax^2$} & {\R{第}{だい}4\R{巻}{かん}から\R{先}{さき}：\bknextpill{5-1}\,\bknextpill{5-2}} \\
{\bknextpill{3-10}\ ボールはいつその\R{高}{たか}さに\R{届}{とど}く？：2\R{次方程式}{じ|ほう|てい|しき}} & {\R{第}{だい}4\R{巻}{かん}から\R{先}{さき}：\bknextpill{5-5}\,\bknextpill{5-7}\,\bknextpill{5-11}\,\bknextpill{6-12}} \\
{\bknextpill{3-11}\ \R{指数関数}{し|すう|かん|すう}：いまある\R{量}{りょう}に\R{比例}{ひ|れい}してふえる・へる} & {\R{第}{だい}4\R{巻}{かん}から\R{先}{さき}：\bknextpill{5-14}\,\bknextpill{5-18}\,\bknextpill{6-4}} \\
{\bknextpill{3-12}\ \R{対数}{たい|すう}：けたを\R{数}{かぞ}える\R{道具}{どう|ぐ}} & {\R{第}{だい}4\R{巻}{かん}から\R{先}{さき}：\bknextpill{5-14}\,\bknextpill{9-8}} \\
{\bknextpill{3-14}\ ばらつきを1つの\R{数}{かず}で：\R{標準偏差}{ひょう|じゅん|へん|さ}} & {\R{第}{だい}4\R{巻}{かん}から\R{先}{さき}：\bknextpill{9-4}} \\
{\bknextpill{3-15}\ 2つの\R{量}{りょう}の\R{関係}{かん|けい}：\R{散布図}{さん|ぷ|ず}と\R{相関}{そう|かん}} & {\R{第}{だい}4\R{巻}{かん}から\R{先}{さき}：\bknextpill{4-12}} \\
{\bknextpill{3-16}\ データに\R{直線}{ちょく|せん}を\R{引}{ひ}く：\R{最小二乗法}{さい|しょう|に|じょう|ほう}} & {\R{第}{だい}4\R{巻}{かん}から\R{先}{さき}：\bknextpill{5-7}\,\bknextpill{7-5}} \\
{\bknextpill{3-17}\ \R{偶然}{ぐう|ぜん}のばらつき：\R{確率}{かく|りつ}と\R{正規分布}{せい|き|ぶん|ぷ}} & {\R{第}{だい}4\R{巻}{かん}から\R{先}{さき}：\bknextpill{5-17}\,\bknextpill{9-4}} \\
\end{longtblr}
